\begin{table*}[t]
\caption{Checklist for a named model and a planned flow, quantity, tolerance and numerical set-up. A missing fact gives \emph{unknown}, not evidence of failure or success. The checklist is a proposal for structuring judgement; its predictive value has not been established.}\label{tab:checklist}
\footnotesize
\begin{tabular*}{\textwidth}{@{\extracolsep{\fill}}p{0.11\textwidth}p{0.31\textwidth}p{0.51\textwidth}@{}}
\toprule
Step & Question & Evidence to record and consequence \\
\midrule
1. Quantity & What is to be predicted, and does the set-up constrain it? & Specify quantity, position, averaging, error measure and tolerance. Derive the relevant conservation constraint. Use a free quantity to assess constitutive accuracy; keep constrained quantities as balance checks. \\
2. Model & Which limitations could matter here? & Record output restrictions, retained terms, available inputs, calibration and tuning, and coupling. Identify relevant flow features and uncertainty. These are candidate explanations, not diagnosed causes. \\
3. Test design & What does the published test actually test? & Distinguish exact-input evaluation from coupled validation. Within coupled tests, distinguish calibration cases from documented transfer cases; unknown calibration remains unknown. Both in-range and out-of-range failures are informative. \\
4. Comparison & What improved, and is the result accurate enough? & Record the matched baseline and reference separately. Compare the same flow, grid, closure, height and spatial support. Keep relative benefit, absolute error, tolerance and uncertainty as separate assessments. No matched baseline means benefit is unassessed; it does not invalidate a reference comparison. \\
5. Planned use & How far can the evidence be extended? & List differences in model version, flow, parameters, quantity and numerics. A direct match supports the stated result within its uncertainty. An extension needs a physical argument and relevant sensitivity evidence at the required tolerance. Otherwise withhold assurance and propose a comparison; do not infer failure. \\
\midrule
Next test & \multicolumn{2}{p{0.85\textwidth}}{Name the uncertainty the test addresses. Specify model and baseline, flow, forcing, grid, closure, matching height, quantities, reference and uncertainty estimates. Score the intended quantity alongside other affected outputs. Use an ablation or controlled change to test a candidate mechanism, recording what still changes with it. State what the test can establish and what remains unresolved. If no reference is available, sensitivity or cross-model agreement alone cannot establish absolute accuracy.} \\
\botrule
\end{tabular*}
\end{table*}
